%! Observational Studies — Unit 1, Lesson 1.5 — Guided Notes & Practice
% THE in-class document, in the MAIN IDEAS / NOTES shape (LESSON_SHAPE.md §1, ported from AP
% Statistics 2026-09-12): the vocabulary box, the hook, then ONE two-column guidednotes table.
% Each row is one idea — a label on the left; on the right one or two COMPLETE printed
% sentences, ONE large pre-drawn display the student reads or names, then a bold prompt and a
% two-across grid of numbered problems with 1.1–1.8cm of writing room. Problems run 1..18.
% DENSITY: no \blank{} at all on this page — every answer is written in a problem's own space
% or in a \labelbox under a display; no mid-sentence blanks; every display is pre-drawn; the
% page belongs to the student's pen.
%   I DO   : the teacher reads the definition, marks up the display, works the first problem
%            of each grid (1, 5, 9, 11).
%   WE DO  : the class works the rest of each grid, students holding the pen; the last row,
%            Guided Practice (15–18), is worked entirely together on Riverbend's Study Center
%            instead of the pool.
%   YOU DO : nothing on this page — ../homework/ IS the individual practice, started in the
%            last ten minutes of the period. No independent set, no reflection box, no
%            objectives box (the cover carries the targets).
% THE TRAP is problem 3 (the director ASSIGNS swim times — the one plan that is not
% observational; students sort by "is it a survey" instead of "who decided the groups").
% THE CRUX is problem 13 (row 4, "Not Causation"): the flyer's claim is not supported even
% though the study has no bias — three stories fit one table, and the hook vote is re-taken
% there (PS.DC.2d). Problem 14 is the prediction-about-switching, a causal claim in disguise.
% Problem 17 is the crux transferred to Riverbend (the memo's arrow runs backwards).
%
% THE DATA (verified in Python before authoring)
%   Harbor Point Community Pool: 1,500 members; SRS of 200 from the full list, 100% response;
%     80 morning swimmers, 56 sleep 7+ hrs -> 56/80 = 70%; 120 evening, 54 -> 54/120 = 45%;
%     gap 25 points; pooled 110/200 = 55%; 0.55 x 1,500 = 825
%     retirement (the lurking variable): 60/80 = 75% of morning vs 18/120 = 15% of evening
%   Guided Practice, Riverbend HS: 900 students; random sample of 150 records; 60 used the
%     Study Center weekly, 39 passed every class -> 65%; 90 never used it, 72 passed -> 80%;
%     gap 15 points, NON-users higher; previous quarter 45/60 = 75% had failed a class vs
%     9/90 = 10% (the arrow runs backwards)
%
% PAGE LOCKSTEP with notes_key/: the key is GENERATED from this file by templates/lesson/mkkey.py
% (spec: templates/lesson/examples/lesson05_notes_key.json) — every \vterm{} becomes a
% \vtermans{}{}, every \par\writelines{n} n \ansline{} calls, and the answer argument of every
% \pcell, \writespace and \labelbox is filled. Answer spaces are fixed heights, so the two
% files cannot drift. KEEP EVERY \pcell ON ONE SOURCE LINE — mkkey matches line by line.
% TABLE MECHANICS: inside a cell, `\\` ENDS THE ROW — break lines with \par. A row cannot break
% across pages: the figure gets its own sub-row (`\\` then `& ...`) and EACH GRID ROW is its own
% sub-row — close the probgrid, `\\ &`, reopen as probgrid* (no top rule).
\documentclass[12pt]{article}
\usepackage{saar-article}
\usepackage{saar-boxes}
\newcommand{\TallMath}[1]{$\displaystyle #1\rule[-1.4em]{0pt}{3.2em}$}
% Word bank strip — ONLY above a table the student fills. This lesson has no fill table, so
% the macro is defined (the key inherits it) but never used.
\newcommand{\wordbank}[1]{\par\vspace{2pt}\noindent\fcolorbox{goldacc}{goldbg}{\parbox{\dimexpr\linewidth-2\fboxsep-2\fboxrule\relax}{\small\textbf{Word bank:} #1}}\par\vspace{4pt}}
% Vocabulary rows: a FIXED-HEIGHT row carrying the term label and OPEN writing space —
% no underline beside the term and no rule line (user correction 2026-09-06: the black
% answer line crowded the row; students write beside and below the term). The key fills
% exactly the same height, so the box cannot drift blank vs. keyed. saar-article's
% \termblank adds an inline blank and a rule line, so the pair is defined here (the key is
% generated from this file and inherits both). 2.3cm per row gives students three
% handwritten lines of room (2.1cm here, so the table starts on page 1); keep key
% definitions to two lines.
\newcommand{\termrowheight}{2.1cm}
\newlength{\termrowinset}\setlength{\termrowinset}{7pt}
\newcommand{\vterm}[1]{%
  \par\noindent\begin{minipage}[t][\termrowheight][t]{\linewidth}%
    \vspace*{\termrowinset}%
    \textbf{\textcolor{cerulean}{#1:}}%
  \end{minipage}\par}
\newcommand{\vtermans}[2]{%
  \par\noindent\begin{minipage}[t][\termrowheight][t]{\linewidth}%
    \vspace*{\termrowinset}%
    \textbf{\textcolor{cerulean}{#1:}}\quad\ans{#2}%
  \end{minipage}\par}

\begin{document}

\pageheader{Unit 1, Lesson 1.5}{Guided Notes \& Practice}
% NAMESTRIP: no \namedateperiod here. The student writes their name once, on the
% cover this component is stapled behind.

\vspace{0.15in}

\begin{vocabbox}
\small
Fill in each term \textbf{as we name it} in the notes below.
\vterm{Observational study}
\vterm{Explanatory variable}
\vterm{Response variable}
\vterm{Association}
\vterm{Lurking variable}
\end{vocabbox}

\boxguard[12]
\begin{hookbox}
\small
Harbor Point's director reads a clean survey --- a random sample of $200$, everyone answered,
every count correct: $70\%$ of morning swimmers get $7$+ hours of sleep; only $45\%$ of evening
swimmers do. Her flyer: \emph{``Swim in the morning --- sleep better!''}

\vspace{3pt}
\textbf{Is the flyer's claim safe?} Circle one: \quad safe \qquad not safe \quad --- and say why
you think so. We vote again at problem~13.
\par\writelines{2}
\end{hookbox}

\small
\begin{guidednotes}
% ── ROW 1: the one word — assign ─────────────────────────────────────────────
\mainidea[One word:]{Assign} &
In an \textbf{observational study} you measure or survey a sample \emph{without trying to
affect anyone}. The researcher \textbf{does not assign} anybody to a group --- the groups
already exist. \textbf{A survey is one type of observational study}, so every study in Unit~1
so far has been one.
\\
& {\centering
\begin{tikzpicture}[scale=0.8]
  % left panel: watching — the members chose their own swim time
  \fill[frost, rounded corners=5pt] (0,0) rectangle (5.4,3.0);
  \draw[cerulean, thick, rounded corners=5pt] (0,0) rectangle (5.4,3.0);
  \node[font=\footnotesize\bfseries\color{cerulean}] at (2.7,2.62) {WATCH};
  \node[font=\scriptsize, align=center] at (2.7,1.85) {Members already swim when\\they like. Ask $200$ of them\\when they swim and how they sleep.};
  \node[font=\scriptsize\bfseries, align=center] at (2.7,0.55) {Who put each member\\in a group?};
  % right panel: assigning — the director puts members into groups
  \fill[goldbg, rounded corners=5pt] (6.2,0) rectangle (11.6,3.0);
  \draw[goldacc, thick, rounded corners=5pt] (6.2,0) rectangle (11.6,3.0);
  \node[font=\footnotesize\bfseries\color{goldacc}] at (8.9,2.62) {ASSIGN};
  \node[font=\scriptsize, align=center] at (8.9,1.85) {The director \emph{puts} $100$ members\\at 6~a.m.\ and $100$ at 7~p.m.\\for a month, then compares sleep.};
  \node[font=\scriptsize\bfseries, align=center] at (8.9,0.55) {Who put each member\\in a group?};
\end{tikzpicture}\par\vspace{5pt}
Left: \labelbox{2.2cm}{}\quad Right: \labelbox{2.2cm}{}\quad Observational: \labelbox{1.5cm}{}\par}
\\
& \notesprompt{Observational study or not? Say \emph{who decided the groups}.}
\begin{probgrid}{|Y|Y|}
\pcell{1}{Ask a random sample of members what time they usually swim and how much they sleep.}{1.1cm}{} & \pcell{2}{Count how many members walk in before $8$~a.m.\ every day for a week.}{1.1cm}{} \\ \hline
\end{probgrid}
\\
& \begin{probgrid}*{|Y|Y|}
\pcell{3}{Assign half the members to swim at $6$~a.m.\ and half at $7$~p.m.\ for a month, then compare their sleep. \textbf{Vote first, then decide.}}{1.1cm}{} & \pcell{4}{Mail the sleep survey to \emph{every} member on the membership list (a census).}{1.1cm}{} \\ \hline
\end{probgrid}
\\ \hline
% ── ROW 2: the statistical cycle plans it (PS.DC.2d) ─────────────────────────
\mainidea[The plan, again:]{The Cycle} &
Lesson~1.1's four stages plan an observational study too --- the same cycle, with Lesson~1.3's
sampling plan in the Collect stage and Lesson~1.4's bias check standing behind it.\par
\vspace{3pt}
\stepnum{1}\ \textbf{Formulate} --- what do we want to know, and about whom?\par
\stepnum{2}\ \textbf{Collect} --- who is on the frame, and how is the sample chosen?\par
\stepnum{3}\ \textbf{Represent} --- how will we organize what comes back?\par
\stepnum{4}\ \textbf{Analyze \& communicate} --- what do the numbers say, \emph{in context}?
\\
& {\centering
\begin{tikzpicture}[scale=0.8]
  % the cycle: four boxes in a loop, Formulate filled in for the pool study
  \fill[frost, rounded corners=5pt] (0,2.0) rectangle (5.4,3.6);
  \draw[cerulean, thick, rounded corners=5pt] (0,2.0) rectangle (5.4,3.6);
  \node[font=\scriptsize\bfseries\color{cerulean}] at (2.7,3.3) {1 \ FORMULATE};
  \node[font=\tiny, align=center] at (2.7,2.6) {Do morning swimmers report\\more sleep than evening\\swimmers?};
  \fill[white, rounded corners=5pt] (6.2,2.0) rectangle (11.6,3.6);
  \draw[cerulean, thick, rounded corners=5pt] (6.2,2.0) rectangle (11.6,3.6);
  \node[font=\scriptsize\bfseries\color{cerulean}] at (8.9,3.3) {2 \ COLLECT};
  \node[font=\tiny\itshape, align=center] at (8.9,2.6) {problem 5};
  \fill[white, rounded corners=5pt] (6.2,0) rectangle (11.6,1.6);
  \draw[cerulean, thick, rounded corners=5pt] (6.2,0) rectangle (11.6,1.6);
  \node[font=\scriptsize\bfseries\color{cerulean}] at (8.9,1.3) {3 \ REPRESENT};
  \node[font=\tiny\itshape, align=center] at (8.9,0.6) {problem 6};
  \fill[white, rounded corners=5pt] (0,0) rectangle (5.4,1.6);
  \draw[cerulean, thick, rounded corners=5pt] (0,0) rectangle (5.4,1.6);
  \node[font=\scriptsize\bfseries\color{cerulean}] at (2.7,1.3) {4 \ ANALYZE};
  \node[font=\tiny\itshape, align=center] at (2.7,0.6) {problem 7};
  % arrows around the loop
  \draw[->, very thick, charcoal] (5.5,2.8) -- (6.1,2.8);
  \draw[->, very thick, charcoal] (8.9,1.9) -- (8.9,1.7);
  \draw[->, very thick, charcoal] (6.1,0.8) -- (5.5,0.8);
  \draw[->, very thick, charcoal] (2.7,1.7) -- (2.7,1.9);
\end{tikzpicture}\par}
\\
& \notesprompt{Fill in what \emph{this} study does at each stage --- then find what the plan never says.}
\begin{probgrid}{|Y|Y|}
\pcell{5}{\textbf{Collect.} Who was on the frame, how was the sample chosen, and how many answered?}{1.2cm}{} & \pcell{6}{\textbf{Represent.} How will the $200$ answers be organized? Name the two things each member is sorted by.}{1.2cm}{} \\ \hline
\end{probgrid}
\\
& \begin{probgrid}*{|Y|Y|}
\pcell{7}{\textbf{Analyze \& communicate.} What do the numbers say, in context? Use your warm-up.}{1.2cm}{} & \pcell{8}{Nowhere in the plan does anybody tell a member when to swim. What does that one absence make the study?}{1.2cm}{} \\ \hline
\end{probgrid}
\\ \hline
% ── ROW 3: two groups that built themselves — explanatory and response ───────
\mainidea[Explanatory \& response]{Two Variables} &
The \textbf{explanatory variable} is the one we think might explain the other --- the tail of
the arrow in the claim. The \textbf{response variable} is the one we measure to see what
happens --- the head of the arrow. Two variables are \textbf{associated} when knowing which
group an individual is in changes what you expect for the other variable.
\\
& {\centering
\begin{tikzpicture}[scale=0.8]
  % the whole study in one two-way table, pre-drawn
  \fill[frost] (0,0) rectangle (8.8,3.3);
  \draw[cerulean, thick] (0,0) rectangle (8.8,3.3);
  \draw[cerulean] (0,2.5) -- (8.8,2.5);
  \draw[cerulean] (0,0.8) -- (8.8,0.8);
  \draw[cerulean] (3.6,0) -- (3.6,3.3);
  \node[font=\scriptsize\bfseries] at (4.9,2.9) {$7$+ hrs};
  \node[font=\scriptsize\bfseries] at (6.6,2.9) {Fewer};
  \node[font=\scriptsize\bfseries] at (8.1,2.9) {Total};
  \node[font=\scriptsize\bfseries, anchor=west] at (0.15,2.05) {Morning swimmers};
  \node[font=\scriptsize] at (4.9,2.05) {$56$}; \node[font=\scriptsize] at (6.6,2.05) {$24$}; \node[font=\scriptsize\bfseries] at (8.1,2.05) {$80$};
  \node[font=\scriptsize\bfseries, anchor=west] at (0.15,1.25) {Evening swimmers};
  \node[font=\scriptsize] at (4.9,1.25) {$54$}; \node[font=\scriptsize] at (6.6,1.25) {$66$}; \node[font=\scriptsize\bfseries] at (8.1,1.25) {$120$};
  \node[font=\scriptsize\bfseries, anchor=west] at (0.15,0.4) {Total};
  \node[font=\scriptsize] at (4.9,0.4) {$110$}; \node[font=\scriptsize] at (6.6,0.4) {$90$}; \node[font=\scriptsize\bfseries] at (8.1,0.4) {$200$};
  % the arrow of the claim
  \draw[->, very thick, goldacc] (9.9,2.6) -- (12.7,2.6);
  \node[font=\scriptsize\bfseries\color{goldacc}, above] at (11.3,2.65) {the claim};
  \node[font=\scriptsize, align=center] at (9.9,1.85) {swim\\time};
  \node[font=\scriptsize, align=center] at (12.7,1.85) {$7$+ hrs\\of sleep?};
  \node[font=\tiny\itshape\color{slate}, align=center] at (11.3,0.7) {nobody built these groups:\\the members chose\\their own swim time};
\end{tikzpicture}\par\vspace{5pt}
Tail of the arrow: \labelbox{2.6cm}{}\qquad Head: \labelbox{2.6cm}{}\par}
\\
& \notesprompt{Read the table. Keep the two denominators straight --- $80$ and $120$.}
\begin{probgrid}{|Y|Y|}
\pcell{9}{From your warm-up: the two percents, the gap in percentage points, and which group is higher.}{1.2cm}{} & \pcell{10}{About how many of all $1{,}500$ members get at least $7$ hours? Show the pooled percent first.}{1.2cm}{} \\ \hline
\end{probgrid}
\\ \hline
% ── ROW 4: THE CRUX — three stories for one table ────────────────────────────
\mainidea[Association is]{Not Causation} &
The $25$-point gap is real, and it is not bias --- the sample was random and everyone
answered. But \textbf{three different stories produce exactly this table.} A variable outside
the study that affects \emph{both} variables in it is a \textbf{lurking variable}.
\\
& {\centering
\begin{tikzpicture}[scale=0.8]
  % story 1: the claim
  \fill[frost, rounded corners=5pt] (0,0) rectangle (3.6,3.4);
  \draw[cerulean, thick, rounded corners=5pt] (0,0) rectangle (3.6,3.4);
  \node[font=\scriptsize\bfseries\color{cerulean}] at (1.8,3.05) {1 \ THE CLAIM};
  \node[font=\scriptsize, align=center] at (1.8,2.2) {morning swim};
  \draw[->, very thick, charcoal] (1.8,1.9) -- (1.8,1.3);
  \node[font=\scriptsize, align=center] at (1.8,1.0) {more sleep};
  % story 2: the arrow backwards
  \fill[frost, rounded corners=5pt] (4.0,0) rectangle (7.6,3.4);
  \draw[cerulean, thick, rounded corners=5pt] (4.0,0) rectangle (7.6,3.4);
  \node[font=\scriptsize\bfseries\color{cerulean}] at (5.8,3.05) {2 \ BACKWARDS};
  \node[font=\scriptsize, align=center] at (5.8,2.2) {more sleep};
  \draw[->, very thick, charcoal] (5.8,1.9) -- (5.8,1.3);
  \node[font=\scriptsize, align=center] at (5.8,1.0) {morning swim};
  \node[font=\tiny\itshape, align=center] at (5.8,0.45) {who can get up\\at 6 a.m.?};
  % story 3: a lurking variable moves both
  \fill[goldbg, rounded corners=5pt] (8.0,0) rectangle (12.8,3.4);
  \draw[goldacc, thick, rounded corners=5pt] (8.0,0) rectangle (12.8,3.4);
  \node[font=\scriptsize\bfseries\color{goldacc}] at (10.4,3.05) {3 \ LURKING};
  \node[font=\scriptsize\bfseries, align=center] at (10.4,2.35) {retired?};
  \draw[->, very thick, charcoal] (10.0,2.05) -- (9.3,1.45);
  \draw[->, very thick, charcoal] (10.8,2.05) -- (11.3,1.45);
  \node[font=\scriptsize, align=center] at (9.3,1.1) {morning\\swim};
  \node[font=\scriptsize, align=center] at (11.4,1.1) {more\\sleep};
  \node[font=\tiny\itshape, align=center] at (10.4,0.35) {retired: $60$ of $80$ morning,\\$18$ of $120$ evening};
\end{tikzpicture}\par\vspace{5pt}
Story 3's ``retired?'' is a \labelbox{3.6cm}{}\par}
\\
& \fcolorbox{cerulean}{frost}{\parbox{\dimexpr\linewidth-2\fboxsep-2\fboxrule\relax}{\centering
\textbf{An observational study can show that two variables are \emph{associated}. It cannot
show that one \emph{causes} the other, because nobody assigned the groups.}}}\par
\vspace{3pt}
Report what you \emph{saw} --- both percents and the gap --- and say the variables are
\textbf{associated}. Never write \emph{causes}, \emph{makes}, \emph{raises}, \emph{improves}.\par
\\
& \notesprompt{Test the flyer three ways, then settle the hook vote.}
\begin{probgrid}{|Y|Y|}
\pcell{11}{\textbf{Story 2.} Write the arrow backwards for \emph{this} table: how could sleep decide swim time, instead of the other way around?}{1.4cm}{} & \pcell{12}{\textbf{Story 3.} Percent retired in each group? Is \emph{retired} inside the table or outside it --- so what is it called?}{1.4cm}{} \\ \hline
\end{probgrid}
\\
& \begin{probgrid}*{|Y|Y|}
\pcell{13}{\textbf{Vote again:} safe \; not safe. Which is supported --- \emph{``morning swimmers report more sleep''} or \emph{``swimming in the morning makes you sleep longer''}? Why only one?}{1.8cm}{} & \pcell{14}{\emph{``If the evening swimmers switched to mornings, they would sleep more.''} Supported? What kind of claim is a prediction about switching?}{1.8cm}{} \\ \hline
\end{probgrid}
\\ \hline
% ── ROW 5: GUIDED PRACTICE — we do, together, a fresh context ────────────────
\mainidea[Guided practice]{Study Center} &
Riverbend High School has $900$ students and an after-school Study Center anyone may walk into.
A counselor samples $150$ records at random. Of the $60$ who used the Center weekly, $39$ passed
every class; of the $90$ who never did, $72$ did. The principal drafts a memo.
\\
& {\centering
\begin{tikzpicture}[scale=0.8]
  \fill[frost] (0,0) rectangle (8.8,3.3);
  \draw[cerulean, thick] (0,0) rectangle (8.8,3.3);
  \draw[cerulean] (0,2.5) -- (8.8,2.5);
  \draw[cerulean] (0,0.8) -- (8.8,0.8);
  \draw[cerulean] (3.4,0) -- (3.4,3.3);
  \node[font=\scriptsize\bfseries] at (4.8,2.9) {Passed all};
  \node[font=\scriptsize\bfseries] at (6.6,2.9) {Did not};
  \node[font=\scriptsize\bfseries] at (8.1,2.9) {Total};
  \node[font=\scriptsize\bfseries, anchor=west] at (0.15,2.05) {Used the Center};
  \node[font=\scriptsize] at (4.8,2.05) {$39$}; \node[font=\scriptsize] at (6.6,2.05) {$21$}; \node[font=\scriptsize\bfseries] at (8.1,2.05) {$60$};
  \node[font=\scriptsize\bfseries, anchor=west] at (0.15,1.25) {Never used it};
  \node[font=\scriptsize] at (4.8,1.25) {$72$}; \node[font=\scriptsize] at (6.6,1.25) {$18$}; \node[font=\scriptsize\bfseries] at (8.1,1.25) {$90$};
  \node[font=\scriptsize\bfseries, anchor=west] at (0.15,0.4) {Total};
  \node[font=\scriptsize] at (4.8,0.4) {$111$}; \node[font=\scriptsize] at (6.6,0.4) {$39$}; \node[font=\scriptsize\bfseries] at (8.1,0.4) {$150$};
  % the memo
  \fill[goldbg, rounded corners=5pt] (9.4,0.6) rectangle (12.9,2.7);
  \draw[goldacc, thick, rounded corners=5pt] (9.4,0.6) rectangle (12.9,2.7);
  \node[font=\scriptsize\bfseries\color{goldacc}] at (11.15,2.35) {THE MEMO};
  \node[font=\scriptsize\itshape, align=center] at (11.15,1.5) {``The Study Center\\is hurting grades.\\Shut it down.''};
\end{tikzpicture}\par}
\\
& \notesprompt{We work these together.}
\begin{probgrid}{|Y|Y|}
\pcell{15}{Observational study? Who decided the groups? Name the explanatory and the response variable.}{1.4cm}{} & \pcell{16}{Compute both pass rates and the gap in percentage points. Which group is higher?}{1.4cm}{} \\ \hline
\end{probgrid}
\\
& \begin{probgrid}*{|Y|Y|}
\pcell{17}{Write the finding as an \emph{association} sentence: both percents, the gap, no causal verb. Does it support the memo?}{1.7cm}{} & \pcell{18}{Last quarter, before the Center opened, $45$ of the $60$ users had already failed a class; $9$ of the $90$ non-users had. Percents? Which story --- and what settles it?}{1.7cm}{} \\ \hline
\end{probgrid}
\\ \hline
\end{guidednotes}

% ── THE NOTES END AT GUIDED PRACTICE ─────────────────────────────────────────
% There is deliberately no "you do" here: ../homework/ is the individual practice,
% started in class in the last ten minutes while the teacher circulates.

\end{document}
